\documentclass{article}

\usepackage[spanish]{babel}
\usepackage{amsmath,amssymb}

\newtheorem{definition}{Definición}

\title{Grafos}
\author{Benjamín Rivas Beltrán}
\date{\today}

\begin{document}

\maketitle

\section{Teoria de grafos}

\begin{definition}{Grafo}

  Un grafo es un par $G = (V, E)$ donde $V$ es un conjunto de vértices y $E$ es un conjunto de aristas, con $E$ siendo un subconjunto de $V \times V$.
\end{definition}

\begin{itemize}
  \item Los vértices son puntos y las aristas son líneas que conectan los vértices.
  \item Un vértice puede tener 0 o más aristas conectadas a él.
  \item Cada una de las aristas conecta exactamente dos vértices.
  \item El orden de un grafo es el número de vértices que tiene. Denotamos el orden de un grafo $G$ como $|V|$.
  \item Si $|V|$ es finito, entonces el grafo se llama grafo finito. 
  \item 
\end{itemize}

\section{Grafos dirigidos y no dirigidos}
\begin{definition}{Grafo no dirigido}

  Un grafo no dirigido es un grafo $G = (V, E)$ donde las aristas no tienen dirección. Es decir, si $(u, v) \in E$, entonces también $(v, u) \in E$.
\end{definition}

\begin{definition}{Grafo dirigido}

  Un grafo dirigido es un grafo $G = (V, E)$ donde las aristas tienen una dirección. Es decir, si $(u, v) \in E$, entonces no necesariamente $(v, u) \in E$.
\end{definition}

\begin{definition}{Extremo de una arista}

Si $(u, v) \in E$, entonces $u$ y $v$ son los extremos de la arista.
\end{definition}
\end{document}